% ICAIF 2026 Submission - ACM sigconf format
% Use: \documentclass[sigconf,anonymous]{acmart} for submission
% Use: \documentclass[sigconf]{acmart} for camera-ready

\documentclass[sigconf,anonymous]{acmart}

% Remove copyright for submission
\setcopyright{none}
\settopmatter{printacmref=true, printfolios=false}
\renewcommand\footnotetextcopyrightpermission[1]{}
\pagestyle{plain}
\acmConference[]{}{}{}

% Packages
\usepackage{booktabs}
\usepackage{algorithm}
\usepackage{algorithmic}
\usepackage{graphicx}

% BibTeX
\bibliographystyle{ACM-Reference-Format}

\begin{document}
\sloppy
\emergencystretch 1.5em

\title{The Rise and Fall of Cross-Factor Predictability: Regime-Conditional HML$\to$SMB Granger Predictability}

\author{Chorok Lee}
\authornote{Corresponding author}
\email{chorok.lee@kaist.ac.kr}
\affiliation{%
  \institution{Korea Advanced Institute of Science and Technology (KAIST)}
  \city{Daejeon}
  \country{South Korea}
}

\begin{abstract}
Cross-factor predictive relationships in equity markets can structurally
break down and not recover.
Using daily Fama-French data (1990--2024) and a Student-$t$ Hidden Markov Model
with $K = 3$ regimes (selected on training data, 1990--2012),
we document that HML Granger-predicts SMB in the pre-crisis Normal regime
($p = 8.75 \times 10^{-9}$, Bonferroni-corrected for 30 factor pairs).
A data-driven sup-$F$ identifies June 1998 as the primary structural break
($p = 1.23 \times 10^{-13}$), with January 2008 confirming further decay
(Chow $p = 2.29 \times 10^{-6}$). The relationship has been
consistent with zero for 16 years (95\% CI $[-0.049, 0.073]$).
This finding is robust across all 7 HMM local optima clusters,
multiple HAC bandwidth specifications, and lags 1--15 in the Normal regime.
A frozen out-of-sample (OOS) test (2013--2024) detects Elevated-regime
patterns (HAC $p = 0.043$, bandwidth-sensitive: NW default $p = 0.056$),
but this does not survive 30-pair Bonferroni correction and is driven by
regime expansion (13.7\%$\to$33.7\%; prevalence-matched bootstrap $p = 0.153$).
The OOS signal exists only at $K = 3$ (null at $K = 2, 4$),
has $\Delta R^2 < 1\%$, and does not generate trading profits
(Sharpe $= -0.07$).
\end{abstract}

\begin{CCSXML}
<ccs2012>
   <concept>
       <concept_id>10002950.10003648.10003671</concept_id>
       <concept_desc>Mathematics of computing~Time series analysis</concept_desc>
       <concept_significance>500</concept_significance>
   </concept>
   <concept>
       <concept_id>10010147.10010257.10010293.10010294</concept_id>
       <concept_desc>Computing methodologies~Causal reasoning and diagnostics</concept_desc>
       <concept_significance>500</concept_significance>
   </concept>
</ccs2012>
\end{CCSXML}

\ccsdesc[500]{Mathematics of computing~Time series analysis}
\ccsdesc[500]{Computing methodologies~Causal reasoning and diagnostics}

\keywords{Factor Investing, Granger Predictability, Regime Switching, Hidden Markov Models, Structural Break, Out-of-Sample Validation}

\maketitle

\section{Introduction}

The August 2007 quantitative meltdown---where systematic equity strategies
lost 27\% in three days~\cite{khandani2011quants}---revealed a blind spot
in factor-based risk management. When leveraged funds held overlapping factor
exposures, forced liquidation by one created cascading price pressure on all
others. Correlation-based risk models could not anticipate this cascade because
they measure co-movement, not temporal precedence: which factor moves first?

\textbf{This paper documents a two-break decay in cross-factor predictability.}
Using daily Fama-French factor returns (1990--2024) and regime-conditional
Granger causality tests, we show that HML$\to$SMB predictive precedence
is regime-specific and undergoes structural breaks at June 1998 and January 2008,
after which the relationship is consistent with zero for 16 years
($p = 0.73$; 95\% CI $[-0.049, 0.073]$).
A frozen out-of-sample design (2013--2024) detects Elevated-regime
patterns, but the signal does not survive multiple-testing correction
and is driven by regime expansion (Section~\ref{sec:frozen_oos}).

\textbf{Methodological scope.} Throughout this paper, ``Granger causality'' refers
to \emph{predictive precedence}---whether past values of one series improve prediction
of another, conditional on the target's own history---not structural causality.
Common drivers (e.g., funding liquidity) cannot be ruled out
(Section~\ref{sec:granger_method}), but the binary appearance and
disappearance of the relationship across regimes rules out static confounding.
We use Student-$t$ Hidden Markov Models to condition Granger tests on latent
market state, discovering regimes unsupervised from distributional properties
and reducing researcher degrees of freedom.

\subsection{Contributions}

\begin{enumerate}
    \item \textbf{Regime-dependent predictive structure with two-break decay.}
    We document that HML$\to$SMB Granger predictability is not merely time-varying
    but \emph{regime-specific}: Bonferroni-significant in the Normal regime
    pre-crisis, absent post-2008. A data-driven sup-$F$ scan identifies
    June 1998 as the primary breakpoint ($p = 1.23 \times 10^{-13}$), with
    January 2008 confirming a second break
    (Chow $p = 2.29 \times 10^{-6}$; Section~\ref{sec:main_result}).
    This challenges the implicit assumption
    in factor-timing models that cross-factor predictive links are
    regime-invariant.

    \item \textbf{Exploratory OOS re-emergence evidence.}
    A frozen OOS design (HMM trained 1990--2012, applied to 2013--2024
    without refitting) detects Elevated-regime re-emergence
    (HAC $p = 0.043$), but this does not survive 30-pair Bonferroni
    correction ($\alpha/30 = 0.00033$). Prevalence sensitivity analysis
    shows the signal is driven by regime expansion
    (13.7\%$\to$33.7\%; bootstrap $p = 0.153$), and it exists only at
    $K = 3$ (null at $K = 2, 4$; Table~\ref{tab:k_sensitivity}).
    This result is valued for its frozen design and economic
    prior rather than statistical significance
    (Sections~\ref{sec:frozen_oos}, \ref{sec:robustness}).
\end{enumerate}

\section{Related Work}

Factor returns exhibit time-series momentum~\cite{ehsani2022factor,gupta2019factor}
and SDF-motivated predictability~\cite{haddad2020factor}, but existing work examines
\emph{within-factor} dynamics---not which factor leads another under which regime.
Chordia and Shivakumar~\cite{chordia2002momentum} and Cooper et al.~\cite{cooper2004market}
establish regime-conditional momentum using observable macro states; our Granger test
conditions on SMB's own lags, isolating incremental \emph{cross-factor} content
(see also~\cite{fama1993common,fama2015five,jegadeesh1993returns,carhart1997persistence,barroso2015momentum,daniel2016momentum,flogel2022momentum,neuhierl2024timing,asness2013value}).

The crowding literature documents that HML/SMB co-exposures amplify crash
risk~\cite{marks2019factor,chincarini2026crowded,betermier2025investor} through
margin spirals~\cite{brunnermeier2009market} and fire
sales~\cite{shleifer1997limits,coval2007asset}
(see also~\cite{stein2009presidential,anton2014connected,lou2022comomentum,kang2021crowding,greenwood2011stock}).
This work measures crowding \emph{intensity} but not regime-conditional
predictive \emph{spillover} between factors.

Hamilton~\cite{hamilton1989new} introduced Markov-switching models;
Bulla~\cite{bulla2011hidden} extends to Student-$t$ HMMs.
Shu and Mulvey~\cite{shu2025dynamic} and Nuriyev et al.~\cite{nuriyev2024regime}
apply regime signals to factor \emph{allocation}, whereas we examine Granger
\emph{precedence between} factors within shared regimes
(see also~\cite{krolzig1997markov,chincoli2017linear,ang2002regime,guidolin2007asset,oelschlager2024fhmm}).

Lo and MacKinlay~\cite{lo1990contrarian} establish cross-autocorrelation between
size-sorted portfolios; Psaradakis et al.~\cite{psaradakis2005markov} pioneer
regime-switching Granger causality; Billio et al.~\cite{billio2012econometric}
and Diebold and Yilmaz~\cite{diebold2012better} develop Granger and FEVD
connectedness measures for financial institutions.
Tank et al.~\cite{tank2022neural} extend Granger causality to nonlinear settings;
we additionally use transfer entropy~\cite{schreiber2000measuring,barnett2009granger}
to characterize complexity.
No prior work examines \emph{regime-conditional} Granger causality
between Fama-French equity \emph{factors}.

\section{Methodology}

\subsection{Data}

We use daily returns for six equity factors from Kenneth French's data library~\cite{french2024data} (five Fama-French factors~\cite{fama2015five} plus Carhart Momentum):
Market excess return (MKT-RF), Size (SMB), Value (HML), Profitability (RMW),
Investment (CMA), and Momentum (MOM)~\cite{jegadeesh1993returns,carhart1997persistence}.

\textbf{Sample:} January 2, 1990 -- December 31, 2024 (8,817 trading days).

\textbf{Scale convention.} We adopt percentage-unit returns (e.g., 0.10 = 0.1\%) as the
\emph{primary} convention. All main-text results, tables, and figures use
percentage units unless explicitly marked ``[decimal-unit]''.
Granger $F$-statistics are scale-invariant; however, HMM emission probabilities are not,
so regime boundaries differ between conventions. Under percentage units,
the frozen OOS yields $n = 953$ usable Elevated-regime days after boundary exclusion; the decimal-unit alternative yields $n = 836$
(regime agreement rate: 86.3\%, with 415 disagreement days concentrated at regime boundaries).
Both yield $F$-$p < 0.05$ in the Elevated regime.
Decimal-unit results appear only as robustness checks in Section~\ref{sec:robustness}.
Daily factor returns are stationary by construction; ADF tests confirm unit-root
rejection ($p < 0.001$) for all six series (BIC lag selection, intercept included,
no trend).

\subsection{Student-\texorpdfstring{$t$}{t} Hidden Markov Model}

Let $z_t \in \{1, \ldots, K\}$ denote the latent regime at time $t$. We model:

\textbf{Transition model:}
$P(z_t = k \mid z_{t-1} = j) = A_{jk}$

\textbf{Emission model (multivariate Student-$t$):}
\begin{equation}
p(\mathbf{x}_t \mid z_t = k) \propto \left(1 + \frac{\delta_k(\mathbf{x}_t)}{\nu_k}\right)^{-\frac{\nu_k+d}{2}}
\end{equation}

where $\delta_k(\mathbf{x}_t) = (\mathbf{x}_t - \boldsymbol{\mu}_k)^\top\boldsymbol{\Sigma}_k^{-1}(\mathbf{x}_t - \boldsymbol{\mu}_k)$.

\textbf{Why Student-$t$?} Financial returns exhibit excess kurtosis.
Student-$t$ distributions accommodate heavy tails, enabling detection of
moderate stress states that Gaussian HMMs miss.
Estimated degrees of freedom ($\hat{\nu}_{\text{Normal}} = 6.2$,
$\hat{\nu}_{\text{Elevated}} = 3.9$, $\hat{\nu}_{\text{Crisis}} = 5.5$)
are all well below the Gaussian limit; HAC-robust Granger tests remain
valid without finite-moment assumptions beyond OLS asymptotic normality.

\textbf{Regime label caveat.} Under the BIC-optimal fit, the ``Crisis'' label
denotes a high-kurtosis statistical state that does \emph{not} align with
calendar crises (0\% of 2008 GFC days assigned to Crisis). We retain the label
for consistency with HMM convention; Section~\ref{sec:local_optima} documents
sensitivity fits that do align with calendar crises.

\textbf{Model selection.} $K = 3$ is pre-specified using BIC
\emph{on training data (1990--2012) alone}
($\Delta\text{BIC}$ over $K = 2$: $1{,}680$; $K = 4$ also disfavored),
eliminating look-ahead bias in regime count selection.
For transparency, full-sample BIC also favors $K = 3$ ($\Delta\text{BIC}$
over $K = 2$: $847$; over $K = 4$: $\Delta\text{BIC} = 312$). Transition
probabilities are estimated via EM with 50 random
initializations.\footnote{Three seeds (28, 20, 6) reach the same full-sample
log-likelihood; seed~28 is primary (sorted-order convention).
Among 16/50 fits (32\%) assigning $\geq$50\% of 2008 days to Crisis,
seed~42 is the highest-LL sensitivity fit (LL cost $<$0.3\%); used for
FF25 overlap (Section~\ref{sec:overlap_evidence}) and VaR
(Appendix~\ref{app:var}), disclosed each time. All tables use the primary
fit unless noted.}

\textbf{Local optima and economic validity.}
Our 50-seed multistart reveals 7 distinct local optima clusters (Section~\ref{sec:local_optima};
Appendix~\ref{app:optima}). The BIC-optimal fit (Cluster~1, seed~28) assigns 0\% of 2008
GFC days to Crisis, while economically sensible fits (Clusters~5--7) assign
90--100\%. We report results under \emph{both} the BIC-optimal (Cluster~1, primary)
and the economically valid (Cluster~5, sensitivity) fits. The $\geq$50\% GFC detection
criterion is a post-hoc sensitivity analysis, not a pre-specified selection rule.
The OOS Elevated raw $F$-$p < 0.05$ in 7/7 clusters; however, this robustness
is conditional on the full OOS sample and does not address the prevalence-dependence
(bootstrap $p = 0.153$).

\subsection{Addressing Regime-Identification Circularity}

In-sample Granger tests condition on regime labels from the same returns,
creating potential circularity. Two mitigations: (1)~the HMM uses
\emph{distributional} properties while Granger tests examine \emph{temporal
dynamics}---distinct statistical features; (2)~soft-label sensitivity (posterior
probabilities rather than Viterbi assignments) yields identical conclusions.

The \emph{primary} validation eliminates regime-identification circularity:
following Welch and Goyal~\cite{welch2008comprehensive}, we fit the HMM on
1990--2012, freeze all parameters, and classify 2013--2024 without refitting
(Section~\ref{sec:frozen_oos}). $K = 3$ is selected on training data alone
($\Delta\text{BIC} = 1{,}680$ over $K = 2$). Pair-selection bias remains
(Section~\ref{sec:pairselection}).

\subsection{Per-Regime Granger Predictability}
\label{sec:granger_method}

Granger causality~\cite{granger1969investigating,shojaie2022granger} tests whether past values
of one series improve prediction of another, conditional on the target's own history.
For each regime $k$, we extract $\mathcal{T}_k = \{t : \hat{z}_t = k\}$
and test:\footnote{We test pairwise (bivariate) Granger relationships, with a
trivariate extension conditioning on MKT-RF (Section~\ref{sec:robustness}).
A full 6-factor VAR ($6^2 \times 9 = 324$ parameters per regime) is
under-identified at $n \approx 1{,}000$; post-double-selection
methods~\cite{hecq2023granger} could address this.}
$H_0: r_{j,t} \perp \{r_{i,t-\ell}\}_{\ell=1}^{L} \mid \{r_{j,t-\ell}\}_{\ell=1}^{L}$

\textbf{Multiple testing framework.}
In-sample analysis (1990--2024) uses a stringent Bonferroni correction:
$\alpha_{\text{fam}} = 0.01$ across 30 directed factor pairs, yielding
$\alpha_{\text{ind}} = 0.00033$. 

For frozen OOS analysis (2013--2024), the appropriate correction depends on the
degree of in-sample screening. Because HML--SMB was identified from screening
all 30 directed pairs in-sample, the \emph{primary} OOS correction is
30-pair Bonferroni ($\alpha_{\text{fam}} = 0.01$, $\alpha/30 = 0.00033$), under
which no OOS pair survives. As a \emph{sensitivity} analysis recognizing
the economic prior for HML--SMB, we also report 3-regime Bonferroni
($\alpha/3 = 0.0167$), under which HAC $p = 0.043$ still does not survive.
The OOS result is valued for its frozen
design and economic motivation, not statistical significance.
HAC standard errors use Andrews~\cite{andrews1991heteroskedasticity}
automatic bandwidth selection (AR(1) plug-in for Bartlett kernel) as the primary specification,
applied consistently across all Granger tests. Bandwidth~1 (matching Granger lag) is a sensitivity check. Full bandwidth sensitivity is in
Section~\ref{sec:robustness}.

\textbf{Regime boundary handling.} For Granger tests at lag $L$,
we require all lags within the same regime ($\hat{z}_{t-\ell} = k$ for
$\ell \in \{1, \ldots, L\}$). Exclusion rates are minimal: 0.67\%
in-sample (59/8{,}817 at lag~1), though 7.4\% for frozen OOS (224/3{,}020)
due to more frequent transitions in 2013--2024. Full exclusion statistics
are in Appendix~\ref{app:optima}.

\subsection{Factor Pair Selection}
\label{sec:pairselection}

We focus on the HML--SMB pair for three reasons. First, \emph{economic
plausibility}: Value and Size factors have documented overlap in institutional
holdings (Section~\ref{sec:interpretation}), providing an economically motivated
mechanism---when leveraged institutions must unwind value positions, the associated
size exposure is simultaneously deleveraged, creating a predictive lag from HML
to SMB. This economic hypothesis precedes and motivates our data analysis.
Second, \emph{in-sample screening}: a preliminary screen of all 30 directed
factor pairs found that HML--SMB shows the most distinctive regime-heterogeneous
pattern, consistent with the economic prior. Crucially, the primary (LL-only) fit
independently identifies HML--SMB as distinctive (Normal: $p = 1.0 \times 10^{-6}$,
Elevated: $p = 0.013$, Crisis: $p = 0.880$ at lag 5)---no 2008-based screening is
involved in pair selection.
Third, \emph{architectural contrast}: lag-1 evidence is strongest in Normal and
substantially weaker outside Normal (below $\alpha/30 = 0.00033$).

\textbf{Pair selection and multiple comparisons.} The OOS frozen test targets
the HML--SMB pair based on an economic prior (liquidity-mediated contagion),
but was not formally pre-registered. Since HML--SMB was identified from screening
30 in-sample pairs, the OOS test inherits this multiplicity burden.
We report results under three correction frameworks:
(i)~\emph{Primary:} 30-pair Bonferroni ($\alpha/30 = 0.00033$): neither HAC $p = 0.043$
nor $F$-$p = 0.003$ survives;
(ii)~Sensitivity: 3-regime Bonferroni ($\alpha/3 = 0.0167$): HAC $p = 0.043$
does not survive, $F$-$p = 0.003$ survives;
(iii)~Uncorrected ($\alpha = 0.05$): both survive.
Additionally, bootstrap prevalence reweighting
yields median $p = 0.153$ (Section~\ref{sec:frozen_oos}).
The OOS Elevated result is valued for its frozen parameters and
economic prior, not statistical significance under any correction framework.
\emph{Pre-registration:} the HML--SMB pair was not pre-registered;
a replication on independent data (e.g., international factor
returns) would provide confirmatory evidence.

\textbf{Permutation test.} A nonparametric permutation
test~\cite{good2005permutation} shuffles HML labels within the Elevated regime
(preserving regime counts and SMB temporal structure) and recomputes the Granger
$F$-statistic. Under the primary percentage-unit convention ($n = 953$,
50{,}000 shuffles): $p = 0.022$. As a robustness check,
the decimal-unit convention ($n = 836$, 10{,}000 shuffles) yields
$p = 0.063$ [decimal-unit].

\textbf{30-pair multiplicity.} We run Granger tests for all 30 directed factor pairs
in the OOS Elevated window (percentage-unit convention, seed~28).
Applying Benjamini-Hochberg FDR correction at $\alpha = 0.05$:
\textbf{no pairs survive} FDR correction on either $F$-$p$ or HAC-$p$.
HML$\to$SMB ranks 2nd of 30 by $F$-statistic
($F = 9.06$, $p = 0.003$); the top-ranked pair is MOM$\to$SMB ($F = 20.3$).
The rank-based max-statistic $p = 2/30 = 0.067$.
The absence of FDR-surviving pairs underscores that the OOS evidence for any
specific factor pair is modest; the economic prior for HML$\to$SMB provides the
primary justification for focused attention on this pair.

\section{Results}

\subsection{Regime Characteristics}

Table~\ref{tab:regimes} summarizes the three identified regimes. Figure~\ref{fig:timeline}
shows regime assignments over the sample period with major crisis events marked.

\begin{table}[t]
\centering
\caption[Regime statistics]{Regime Summary Statistics (1990--2024, full-sample primary fit). Transition probabilities are estimated via Expectation-Maximization. The frozen OOS analysis (Section~\ref{sec:frozen_oos}) uses a separate fit trained on 1990--2012 only, yielding different regime counts. Elevated and Crisis regimes have similar mean norms (2.03 vs.\ 2.09) but are distinguished by tail shape ($\nu = 3.9$ vs.\ $5.5$) and covariance structure.}
\label{tab:regimes}
\small
\setlength{\tabcolsep}{4pt}
\begin{tabular}{lccccc}
\toprule
Regime & Days & Prop. & Mean $\|\mathbf{x}\|$ (daily \%) & $\nu$ & $P(z_t=z_{t-1})$ \\
\midrule
Normal & 4,723 & 53.6\% & 0.98 & 6.2 & 0.994 \\
Elevated & 3,023 & 34.3\% & 2.03 & 3.9 & 0.991 \\
Crisis & 1,071 & 12.1\% & 2.09 & 5.5 & 0.993 \\
\bottomrule
\end{tabular}
\setlength{\tabcolsep}{6pt}
\end{table}

\begin{figure}[t]
\centering
\includegraphics[width=\columnwidth]{figures/regime_timeline.pdf}
\caption{Regime assignments over sample period (1990--2024). Top panel shows daily
factor volatility (6-factor norm) with regime-colored background. Bottom panel shows
regime classifications. Vertical dashed lines mark major stress events (LTCM/Russia 1998, dot-com 2000--02,
GFC 2008, COVID-2020, 2022 rate shock). Under the
primary (leakage-safe) fit, Crisis regime (red) reflects a high-kurtosis statistical
state; the 2008-screened sensitivity fit aligns it with calendar crises (see
Section~\ref{sec:gaussian_comparison}). Elevated regime (yellow) captures moderate
volatility periods.}
\Description{Two-panel timeline showing daily factor-volatility norm with background regime colors and event markers, and a lower panel with daily regime assignments from 1990 to 2024.}
\label{fig:timeline}
\end{figure}

\textbf{Regime label note.}\label{sec:gaussian_comparison} Under the leakage-safe primary fit, the ``Crisis''
regime captures a high-kurtosis statistical state that does not align with
calendar crises (0\% of 2008, 2011, 2020 days assigned to Crisis vs.\ 83\%
under a Gaussian baseline).\footnote{Table comparing Student-$t$ and Gaussian
crisis detection rates: 2008 Financial (0\% vs.\ 83.3\%), 2011 EU Debt
(0\% vs.\ 25.9\%), 2020 COVID-19 (0\% vs.\ 81.7\%).}
We retain the label for consistency with the HMM
literature~\cite{bulla2011hidden}; the structural break finding is time-indexed
(January 2008), not regime-indexed, and does not depend on this assignment.
The 2008-screened sensitivity fit (seed~42) does align with calendar crises;
we report it separately where noted.

\subsection{The Structural Break in HML$\to$SMB Predictability}
\label{sec:main_result}

\begin{table}[t]
\centering
\caption{Granger Predictability Between HML and SMB by Regime (1990--2024, in-sample,
BIC-optimal fit, seed~28). Lag selected by BIC. Significance threshold:
$p < 0.00033$ (Bonferroni-corrected, $\alpha_{\text{fam}} = 0.01$, 30 directed factor pairs).
Only Normal-regime HML$\to$SMB survives correction.}
\label{tab:main}
\small
\begin{tabular}{llccc}
\toprule
Regime & Direction & $p$-value & Lag & Significant \\
\midrule
Normal & HML $\to$ SMB & $\mathbf{8.75 \times 10^{-9}}$ & 1 & \textbf{Yes} \\
Normal & SMB $\to$ HML & $0.864$ & 1 & No \\
\midrule
Elevated & HML $\to$ SMB & $0.004$ & 1 & Suggestive$^*$ \\
Elevated & SMB $\to$ HML & $0.897$ & 1 & No \\
\midrule
Crisis & HML $\to$ SMB & $0.695$ & 1 & No \\
Crisis & SMB $\to$ HML & $0.294$ & 1 & No \\
\bottomrule
\multicolumn{5}{l}{\footnotesize $^*$Below 0.01 but above Bonferroni threshold $\alpha/30 = 0.00033$.}
\end{tabular}
\end{table}

\textbf{Fit note.} Table~\ref{tab:main} uses the BIC-optimal fit (Cluster~1, seed~28; Section~\ref{sec:local_optima}). The structural break finding is time-indexed (June 1998, January 2008) and robust across all 7 local optima clusters.

\textbf{Key finding under Bonferroni-significant threshold:} only Normal-regime
HML$\to$SMB survives Bonferroni correction.
\begin{itemize}
    \item \textbf{Normal (Bonferroni-significant):} Strong lag-1 effect
    ($p = 8.75 \times 10^{-9}$), robust to HAC correction. The Granger test conditions
    on lagged SMB, so the HML coefficient reflects \emph{incremental cross-factor}
    predictive content beyond SMB's own history---not a within-SMB autocorrelation artifact.
    \item \textbf{Elevated (suggestive):} HML$\to$SMB is directional at
    conventional levels ($p = 0.004$) but does not pass the paper's
    multiple-testing threshold.
    \item \textbf{Crisis (null):} No detectable lag-1 Granger effect in either
    direction.
\end{itemize}

\textbf{Robustness to heteroskedasticity.}
Financial returns exhibit volatility clustering, which may inflate standard Granger
$F$-test significance. We recompute at lag 1 (BIC-optimal) using Newey--West~\cite{newey1987simple}
heteroskedasticity-and-autocorrelation-consistent (HAC) standard errors.
Normal-regime HML$\to$SMB survives HAC correction
($p = 9.45 \times 10^{-8}$ HAC vs.\ $p = 8.75 \times 10^{-9}$ standard).
Elevated HML$\to$SMB weakens under HAC (1990--2024 in-sample: $p = 0.034$), and Crisis remains
non-significant ($p = 0.711$ HAC vs.\ $p = 0.695$ standard).

\textbf{Temporal stability and two-break decay.}
Pre-2008 Normal ($n = 3{,}140$): $p = 6.66 \times 10^{-16}$ ($\Delta R^2 = 2.06\%$).
Post-2008 Normal ($n = 1{,}557$): $p = 0.73$ ($\Delta R^2 < 0.01\%$).\footnote{The
sum $3{,}140 + 1{,}557 = 4{,}697$ is 26 days fewer than Table~\ref{tab:regimes}'s
Normal total of $4{,}723$, due to regime-boundary exclusion at lag~1.}

\textbf{Primary structural break: June 1998.}
The Quandt-Andrews sup-$F = 21.2$ maximizes at \textbf{June 1998} (LTCM/Russia crisis):
$F = 21.20$, $p = 1.23 \times 10^{-13}$.
The top-5 break candidates all cluster in 1998--2003
(June 1998, July 1998, April 1998, August 2003, March 1998),
suggesting the initial weakening of HML$\to$SMB predictability began with
LTCM-driven liquidity stress rather than the GFC.
This timing is economically interpretable: the LTCM collapse triggered
the first major cross-factor deleveraging event in modern quantitative finance,
permanently altering the informational structure of factor markets.

\textbf{Secondary break: January 2008.}
A theory-motivated Chow test at January 2008 (testing stability of the
intercept, HML lag, and SMB lag coefficients) confirms further deterioration:
$F(3,n-6) = 9.68$, $p = 2.29 \times 10^{-6}$; $\hat{\beta}_{\text{HML}}$ shifts
from $-0.189$ (pre-GFC) to $+0.010$ (post-GFC, Wald $z = 5.05$, $p = 9.2 \times 10^{-7}$).
This is $10^{7}\times$ weaker than the June 1998 break and should be interpreted as
confirming a \emph{later phase} of two-break structure rather than identifying the
primary breakpoint. Together, the evidence supports gradual erosion of HML$\to$SMB
predictability beginning around June 1998, not a single GFC-triggered collapse.
A formal Bai-Perron~\cite{bai2003computation} sequential multiple-break test
would provide confirmatory evidence for the number of breaks; we note the
two-break characterization is supported by the sup-$F$ scan and the
theory-motivated Chow test but not formally tested via sequential procedures.

\textbf{Post-break coefficient magnitude.}
The post-2008 coefficient ($\hat{\beta} = 0.012$, SE $= 0.031$, 95\% CI $[-0.049, 0.073]$)
is small relative to the pre-break value ($|\beta_{\text{pre}}| = 0.189$).
The 95\% CI is the primary measure of post-break effect magnitude.
As a sensitivity analysis, a TOST equivalence test~\cite{schuirmann1987comparison}
yields margin-dependent results: at $\delta = 0.095$ (50\% of pre-break effect),
$p_{\text{TOST}} = 0.004$; at $\delta = 0.05$, $p_{\text{TOST}} = 0.110$
(not significant). Because the equivalence margin was not pre-specified,
the TOST is a sensitivity check, not a confirmatory test.
Post-GFC attenuation may reflect regulatory constraints
(Volcker Rule reducing proprietary factor trading), increased factor crowding
(HML--SMB correlation rose from $\rho_{\text{pre}} = 0.01$ to $\rho_{\text{post}} = 0.23$,
potentially absorbing the lagged predictive link into contemporaneous co-movement),
or two-break structure (sup-$F$ at June 1998 hints the single-break Chow test
may not fully resolve the regime transition timing).
These are post-hoc interpretations; the break itself is the finding.

\textbf{Economic magnitude.}
Adding HML lags (1--15) to an SMB autoregression increases $R^2$ by
$\Delta R^2 = 1.04\%$ in Normal ($p = 5.98 \times 10^{-5}$),
1.32\% in Elevated ($p = 0.002$), and 0.56\% in Crisis ($p = 0.988$).
Pre-GFC Normal $\Delta R^2 = 2.06\%$---modest but typical for daily
factor predictability~\cite{ehsani2022factor}. These magnitudes confirm
statistically significant predictive content without implying
trading profitability (see Section~\ref{sec:limitations}).

\begin{figure}[t]
\centering
\includegraphics[width=\columnwidth]{figures/granger_heatmap.pdf}
\caption{All-pairs Granger causality by regime ($-\log_{10}(p)$, lag=5).
Bordered cells with asterisks survive Bonferroni correction.
The strongest regime-differential pair is MKT-RF$\to$SMB; HML$\to$SMB appears as
a secondary differential pattern.}
\Description{Three heatmaps of pairwise Granger significance by regime; darker cells indicate lower p-values, with Bonferroni-surviving cells marked.}
\label{fig:heatmap}
\end{figure}

Figure~\ref{fig:heatmap} shows Granger causality $p$-values for all 30
directed factor pairs across three regimes (lag~5). HML$\to$SMB is one of
several regime-sensitive pairs; its distinctive feature is the structural
break documented above.

\subsection{Economic Interpretation and Supporting Evidence}
\label{sec:interpretation}\label{sec:overlap_evidence}
We hypothesize a deleveraging cascade channel: forced unwinding of value positions
simultaneously deleverages associated size exposures, creating a predictive lag from
HML to SMB (Appendix~\ref{app:mechanism}). FF25 portfolio overlap analysis (sensitivity
fit, seed~42) finds Granger significance concentrating in small-cap portfolios
($\rho_s = 0.35$, permutation $p = 0.046$; Appendix~\ref{app:overlap}), consistent
with this channel. Event-level validation across six stress episodes shows the expected
pattern in 2 of 6 events, with 4 of 6 directionally correct---heterogeneous and low-power
(Appendix~\ref{app:events}).
Alternative explanations (common drivers, information propagation) cannot be ruled out.

\subsection{Machine Learning Validation and Nonlinear Asymmetry}
\label{sec:ml_validation}

We apply multiple model classes following the
Tank et al.~\cite{tank2022neural} neural Granger protocol to characterize
the \emph{complexity boundary} of cross-factor predictability.
This serves two purposes: (1)~functional-form robustness---if nonlinear
methods detect the same causal direction, the finding is model-agnostic;
(2)~\emph{nonlinear causal discovery}---transfer entropy can reveal
directional information flows invisible to linear tests, providing
a methodological contribution beyond standard Granger analysis.

\textbf{LSTM attention reveals regime-conditional lag structure.}
An attention-mechanism LSTM~\cite{tank2022neural} (32 hidden units,
temporal attention over 9 lags) trained per regime exposes which lag
distances carry predictive content (Figure~\ref{fig:ml_diagnostic},
top row). In the Normal regime, 68.2\% of attention concentrates on
lag~1---a sharp peak consistent with the lag-1 BIC selection
(Table~\ref{tab:main}). This concentration decays monotonically across
regimes: 52.9\% in Elevated, 44.2\% in Crisis (approaching the uniform
baseline of 11.1\%), mirroring the two-break decay of the Granger
signal. The attention gradient provides a model-based consistency check that
the structural break is mechanism-level, not an artifact of test specification.

\textbf{Random Forest importance confirms linear lag dominance.}
A four-model diagnostic (OLS, RF with 100 trees, MLP with 64-32 architecture,
LSTM; full details in Appendix~\ref{app:robustness}) finds that
\emph{no nonlinear method improves prediction in any regime under the primary fit (seed~28)}
(Table~\ref{tab:neural}; RF $p > 0.13$, MLP $p > 0.20$, LSTM $p > 0.63$
under the primary fit). Under an alternative regime partition (sensitivity fit, seed~42), RF shows significant improvement in the Elevated regime ($p = 0.010$) and Crisis ($p = 0.005$), indicating that the ``predominantly linear'' characterization is fit-dependent.
RF permutation importance
(Figure~\ref{fig:ml_diagnostic}, bottom row) shows HML lag-1 as the
dominant feature in the Crisis regime (importance $= 0.043$,
$4\times$ the mean), while Normal-regime HML lags are diffuse---consistent
with the linear Granger test capturing all available predictive content.

\textbf{Transfer entropy reveals nonlinear reverse channel.}
While linear methods detect only forward HML$\to$SMB, transfer
entropy~\cite{schreiber2000measuring,barnett2009granger} (Frenzel--Pompe
kNN, $k = 5$, 200 permutations; Figure~\ref{fig:te_asymmetry})
reveals the \emph{reverse} channel SMB$\to$HML is nonlinear and
substantially stronger in the Normal regime ($z = 5.37$, $p < 10^{-6}$
vs.\ forward $z = 2.45$, $p = 0.007$; Table~\ref{tab:te}).
This asymmetry---linear forward, nonlinear reverse---reconciles the null
reverse Granger ($p = 0.864$) with highly significant reverse TE:
the reverse channel is real but invisible to linear tests.
Both directions collapse in the Crisis regime ($z < 1.3$, $p > 0.11$),
confirming the structural break destroys all predictive structure.

\textbf{Directional asymmetry as a structural finding.}
The transfer entropy asymmetry---linear forward, nonlinear reverse---is a
substantive finding. Standard linear Granger tests capture the full
HML$\to$SMB channel but miss the stronger SMB$\to$HML reverse channel.
Linear VAR-based connectedness measures~\cite{diebold2012better} would
systematically underestimate this information flow.
For the ICAIF community, this suggests that neural Granger
methods~\cite{tank2022neural} may detect cross-factor structure invisible
to classical approaches, motivating regime-conditional nonlinear
factor-timing models as future work.

\begin{figure}[t]
\centering
\includegraphics[width=\columnwidth]{figures/ml_attention_importance.pdf}
\caption{ML diagnostics by regime. \emph{Top:} LSTM attention weights
over 9 lags (dashed line $= 1/9$ uniform baseline). Lag-1 concentration
decays from 68.2\% (Normal) to 44.2\% (Crisis), mirroring the Granger
structural break. \emph{Bottom:} Random Forest permutation importance
for SMB lags (blue) vs.\ HML lags (red). HML lag-1 dominates in Crisis;
Normal-regime HML importance is diffuse.}
\Description{Six-panel bar chart: top row shows LSTM attention weights by lag for Normal, Elevated, and Crisis regimes; bottom row shows RF feature importance for SMB and HML lags by regime.}
\label{fig:ml_diagnostic}
\end{figure}

\begin{figure}[t]
\centering
\includegraphics[width=0.85\columnwidth]{figures/te_asymmetry_by_regime.pdf}
\caption{Transfer entropy directional asymmetry. Forward HML$\to$SMB
(blue) is modest; reverse SMB$\to$HML (red) dominates in Normal
($z = 5.37^{***}$). Both collapse in Crisis. Stars denote
permutation significance ($^{*}p < 0.05$, $^{**}p < 0.01$,
$^{***}p < 0.001$).}
\Description{Grouped bar chart of transfer entropy z-scores by regime and direction, showing reverse SMB-to-HML dominance in Normal regime.}
\label{fig:te_asymmetry}
\end{figure}

\subsection{Frozen-Parameter Out-of-Sample Test (Primary Validation)}
\label{sec:frozen_oos}

Two-stage approaches risk circularity: the HMM uses the same returns for
regime assignment that Granger tests subsequently analyze.
We resolve this by fitting the Student-$t$ HMM exclusively on 1990--2012
($n = 5{,}797$ days), freezing all parameters
($\boldsymbol{\mu}_k$, $\boldsymbol{\Sigma}_k$, $\nu_k$, $\mathbf{A}$),
and classifying 2013--2024 ($n = 3{,}020$ days) without refitting.
Granger tests are then run on the held-out labels using the same
\texttt{select\_lag\_bic} + \texttt{run\_granger\_at\_lag} pipeline as Table~\ref{tab:main},
ensuring full methodological consistency.

Under the primary percentage-unit convention, the frozen HMM assigns
25.6\% to Normal ($n = 773$), 33.7\% to Elevated
($n = 1{,}018$), and 40.7\% to Crisis ($n = 1{,}229$)
in the held-out period (vs.\ 21.5\%, 13.7\%, 64.8\% in training).
The Elevated-regime share more than doubles from 13.7\% (training)
to 33.7\% (test).
After lag-1 boundary exclusion (same rule as Section~\ref{sec:granger_method}),
the OOS retains $724 + 953 + 1{,}119 = 2{,}796$ of $3{,}020$ days (7.4\% excluded).
The OOS exclusion rate (7.4\%) is higher than the full-sample lag-1 Granger exclusion
(0.67\%; 59/8{,}817), reflecting more frequent regime transitions when the
frozen 1990--2012 classifier is applied to the unseen 2013--2024 period.

\begin{table}[t]
\centering
\setlength{\tabcolsep}{5pt}
\caption{Frozen OOS Granger Predictability (HML$\to$SMB), 2013--2024.
HMM trained on 1990--2012 only; relabeling uses train-period ordering;
lag $= 1$ pre-fixed from in-sample BIC; filtered probabilities (no future info).
HAC-$p$ uses Andrews (1991) automatic bandwidth selection (AR(1) plug-in, bandwidth=2);
NW default (bandwidth=6) yields $p = 0.056$.
Bootstrap $p$ = median from 1,000 subsamples reweighted to training-period
Elevated prevalence (13.7\%); prevalence sensitivity shows the raw signal
is driven by regime expansion (13.7\%$\to$33.7\%).
No OOS pair survives 30-pair Bonferroni ($\alpha/30 = 0.00033$).}
\label{tab:frozen_granger}
\small
\begin{tabular}{@{}lccccc@{}}
\toprule
Regime & $n$ & $F$-$p$ & HAC-$p$ & Bootstrap $p$ & $\Delta R^2$ \\
\midrule
Normal    & 724 & 0.185 & 0.190 & 0.188 & 0.25\% \\
Elevated  & 953 & \textbf{0.003} & \textbf{0.043} & \textbf{0.153} & \textbf{0.73\%} \\
Crisis    & 1,119 & 0.314 & 0.328 & 0.331 & 0.07\% \\
\bottomrule
\end{tabular}
\end{table}

\textbf{Result} (Table~\ref{tab:frozen_granger}). The OOS Elevated regime yields
$F$-$p = 0.003$, HAC $p = 0.043$ (Andrews automatic bandwidth),
$\Delta R^2 = 0.73\%$ at lag~1. Normal ($p = 0.185$) and Crisis ($p = 0.314$) are null.
Under 30-pair Bonferroni correction ($\alpha/30 = 0.00033$), the OOS result does not
survive; under 3-regime Bonferroni ($\alpha/3 = 0.0167$), the HAC $p = 0.043$ also fails.

\textbf{Prevalence sensitivity:} bootstrap reweighting to match training-period
Elevated prevalence (13.7\% vs.\ 33.7\% in test) yields median $p = 0.153$,
with only 9.9\% of subsamples significant at $p < 0.05$.
The Elevated regime's post-GFC expansion is the primary driver of significance;
at training-period prevalence, the OOS signal is null.
\emph{Convention note:} Table~\ref{tab:frozen_granger} uses the percentage-unit
pipeline ($n = 953$), consistent with the primary convention.
Decimal-unit results ($n = 836$) are qualitatively similar and reported in
Section~\ref{sec:robustness}.
This result does not survive Bonferroni correction across 30 pairs
($\alpha/30 = 0.00033$); its evidential weight
rests on the economic motivation for the HML--SMB pair
(Section~\ref{sec:pairselection}).
A full 50-seed sensitivity analysis across all 7 local optima clusters finds that
\textbf{all 7 clusters} show Elevated significance ($F$-$p < 0.05$ and HAC-$p < 0.05$
in each), including the highest-BIC cluster (Cluster~7, 6 seeds, BIC$= 76{,}137$;
Elevated HAC-$p = 0.026$).
This represents robustness across 7 distinct HMM solutions spanning the full
range of local optima (BIC 75{,}587--76{,}137).
All results use forward-only (filtered) probabilities and train-period relabeling,
eliminating look-ahead.

\textbf{Power analysis.} At $n = 953$ (percentage-unit primary convention),
the OOS test has $>$99\% power to detect the in-sample effect size
($\Delta R^2 = 2.06\%$) and $\approx$75\% at the observed OOS magnitude
($\Delta R^2 = 0.73\%$); the decimal-unit alternative ($n = 836$) yields
98.7\% and 69.7\%, respectively. The smaller OOS effect reflects genuine
attenuation, not low power.

\textbf{Regime redistribution and distributional stability.}
Elevated grows from 13.7\% (training) to 33.7\% (test), raising the question
of whether the OOS Elevated regime has different distributional properties.
The training-period Elevated has mean factor norm $\|\mathbf{x}\| = 2.01$\% and
$\hat{\nu} = 3.9$; the OOS Elevated has $\|\mathbf{x}\| = 1.87$\% and similar
tail behavior. The moderate difference in means suggests the frozen classifier
assigns some formerly Normal observations to Elevated in the post-GFC period,
consistent with a secular increase in factor volatility rather than model
misspecification. Nonetheless, this distributional shift partially drives the OOS
significance and is a limitation.

\textbf{Prevalence reweighting.}
Bootstrap subsampling (1{,}000 iterations, target $n = 697$ matching training
prevalence) yields median $p = 0.153$; only 9.9\% of subsamples achieve $p < 0.05$.
The OOS significance is partially driven by regime expansion.

\textbf{OOS trivariate control.}
Conditioning on MKT-RF in the OOS Elevated regime, the HML$\to$SMB signal retains
similar magnitude (HAC $p = 0.051$), while MKT-RF contributes no incremental content
($F$-$p > 0.35$), consistent with the in-sample trivariate result.

Event-level OOS behavior is heterogeneous across held-out stress episodes
(Appendix~\ref{app:events}), highlighting that the per-regime pooled result
captures the aggregate signal more reliably than individual events.

\section{Discussion}

\subsection{The Structural Break: Robustness of the Main Finding}
\label{sec:robustness}

The pre-crisis Normal-regime finding ($p = 8.75 \times 10^{-9}$, two-break
confirmed at June 1998 and January 2008) is robust across multiple dimensions.
HAC bandwidth sensitivity (Table~\ref{tab:bandwidth_sensitivity}): the
Andrews~\cite{andrews1991heteroskedasticity} automatic bandwidth
(AR(1) plug-in for Bartlett kernel) yields bandwidth $= 2$
(residual $\hat{\rho} = -0.022$; HAC $p = 0.043$);
the Newey--West default bandwidth $= 6$ yields $p = 0.056$
(\emph{crossing the 5\% threshold}); the rule-of-thumb bandwidth $= 7$
yields $p = 0.064$.
In other words, the OOS HAC significance depends on bandwidth choice:
under Andrews automatic selection the result is marginal ($p = 0.043$),
while under the Newey--West default it is not ($p = 0.056$).
At bandwidth $\leq 4$ the result holds ($p \leq 0.048$); at bandwidth
$\geq 6$ it weakens ($p \geq 0.056$). This bandwidth sensitivity
reinforces the fragility of the OOS finding.

\begin{table}[t]
\centering
\caption{HAC Bandwidth Sensitivity: OOS Elevated HML$\to$SMB $p$-values.
Primary specification: Andrews (1991) automatic AR(1) plug-in (bandwidth~$B = 2$).}
\label{tab:bandwidth_sensitivity}
\small
\begin{tabular}{@{}lc@{}}
\toprule
Bandwidth & HAC $p$-value \\
\midrule
$B = 1$ (Granger lag) & 0.041 \\
$B = 2$ (Andrews automatic) \textbf{[PRIMARY]} & \textbf{0.043} \\
$B = 4$ & 0.048 \\
$B = 6$ (NW default) & 0.056 \\
$B = 10$ & 0.078 \\
$B = 15$ & 0.120 \\
$B = 20$ & 0.148 \\
$B = 30$ & 0.173 \\
\bottomrule
\end{tabular}
\end{table}

\textbf{Trivariate confounder control.}
A common-driver concern is that MKT-RF moves both HML and SMB with different lags,
creating spurious HML$\to$SMB predictability.
Conditioning on one lag of MKT-RF in a trivariate Granger specification, the
Normal-regime HML$\to$SMB signal remains Bonferroni-significant, and
MKT-RF itself contributes no incremental predictive content ($F$-$p > 0.43$).
This does not rule out all common drivers but addresses the most prominent candidate.
BIC selects lag~1 for all regime$\times$direction combinations; significance
is robust across lags 1--15 in Normal (Figure~\ref{fig:lag_sensitivity}).
Rolling 3-year unconditional Granger tests (Figure~\ref{fig:rolling}) show
episodic peaks during stress periods (2008--2009, 2020), consistent with
the regime-conditional finding. A four-model diagnostic (OLS, RF, MLP, LSTM;
Appendix~\ref{app:robustness}) finds no nonlinear improvement for forward
HML$\to$SMB, confirming the linear Granger test captures the full predictive
content. The post-2008 absence ($p = 0.73$) is scope definition, not a
limitation: it demonstrates a clean structural break.

\begin{figure}[t]
\centering
\includegraphics[width=\columnwidth]{figures/rolling_granger.pdf}
\caption{Rolling 3-year unconditional Granger causality ($-\log_{10}(p)$, HML$\to$SMB,
lag=9) with regime-colored background.}
\Description{Time-series plot of rolling Granger significance with shaded regime bands, showing time variation in HML-to-SMB predictability.}
\label{fig:rolling}
\end{figure}

\begin{figure}[t]
\centering
\includegraphics[width=\columnwidth]{figures/lag_sensitivity.pdf}
\caption{Granger causality $p$-values for HML$\to$SMB across lags 1--15 by regime.}
\Description{Line chart of HML-to-SMB p-values by lag for each regime, highlighting strong Normal-regime significance and weak Crisis-regime evidence.}
\label{fig:lag_sensitivity}
\end{figure}

\subsection{OOS Re-Emergence: Exploratory Evidence with Disclosed Fragility}

As detailed in Section~\ref{sec:frozen_oos}, the frozen OOS Elevated result
is sensitive to prevalence (bootstrap $p = 0.153$),
scale convention (permutation $p \in [0.022, 0.063]$),
and regime count (null at $K = 2, 4$; Table~\ref{tab:k_sensitivity}),
though robust across all 7 local optima clusters (Appendix~\ref{app:optima}).
It is informative because it uses frozen parameters and held-out data,
but should not be treated as confirmatory.

\begin{table}[t]
\centering
\caption{$K$-Sensitivity: Frozen OOS HML$\to$SMB by Number of Regimes.
Permutation: 5{,}000 label shuffles (percentage-unit convention).
``Best Regime'' = regime with lowest HML$\to$SMB $p$-value in each specification.}
\label{tab:k_sensitivity}
\scriptsize
\begin{tabular}{@{}lclcccc@{}}
\toprule
$K$ & $\Delta$BIC$_{K\text{vs.}3}$ & Best Regime & $n$ & $F$-$p$ & Perm.\ $p$ & Sig.? \\
\midrule
2 & $+1{,}680$ & High-Vol & 1{,}917 & 0.514 & 0.572 & No \\
\textbf{3} & \textbf{---} & \textbf{Elevated} & \textbf{953} & \textbf{0.003} & \textbf{0.025} & \textbf{Yes} \\
4 & $+312$ & Low-Elevated & 855 & 0.056 & --- & No \\
\bottomrule
\end{tabular}
\end{table}

The non-monotonic pattern---null at $K = 2$, significant at $K = 3$, null at $K = 4$---indicates that the OOS Elevated result depends on the three-regime specification. While BIC strongly favors $K = 3$ ($\Delta\text{BIC} = 1{,}680$ over $K = 2$), this model-structure sensitivity reinforces the fragility of the OOS finding.

\textbf{Pair selection transparency.}
MOM$\to$SMB---the top-ranked pair by OOS $F$-statistic (rank 1/30,
$F = 20.3$)---shows the same regime-conditional pattern (permutation
$p = 0.010$, $\Delta R^2 = 1.94\%$). Our primary pair selection is guided
by an economic prior (HML--SMB institutional crowding) rather than pure
empirical strength. This reduces degrees of freedom but introduces potential
prior-driven selection; pre-registered replication on international data would
provide definitive confirmation.
Additional robustness (subsample stability, MS-VAR, weekly aggregation,
transfer entropy, filtered vs.\ smoothed) is in Appendix~\ref{app:robustness}.

\subsection{Local Optima and Regime Definition Stability}
\label{sec:local_optima}

A critical tension emerges when we compare statistical fit to economic interpretability.
Student-$t$ HMM estimation via EM is sensitive to random initialization; our 50-seed multistart
reveals 7 distinct local optima clusters (Table~\ref{tab:local_optima}) with substantially different crisis-regime assignments.
The BIC-optimal fit (Cluster~1, 3 seeds, BIC$=75{,}587$) assigns \emph{zero percent} of 2008
crisis days to the Crisis regime ($p = 0.88$), yet economically sensible fits (Clusters~5--7)
assign 90--100\% of the GFC to Crisis.

\begin{table}[t]
\centering
\caption{Local Optima Taxonomy: 7 Clusters from 50-Seed HMM Multistart.
BIC-optimal is Cluster~1 (lowest BIC); economically valid clusters ($\geq 50\%$ GFC detection)
are Clusters~5--7. ``\% Elev.\ OOS'' = Elevated regime prevalence in the OOS period (2013--2024);
training-period prevalence is 13.7\% for all clusters. Elevated OOS HAC-$p < 0.05$ in all clusters.}
\label{tab:local_optima}
\scriptsize
\begin{tabular}{@{}lcccccc@{}}
\toprule
Cluster & Seeds & BIC & \% GFC in Crisis & \% Elev.\ OOS & Elev.\ HAC-$p$ & Crisis Perm.\ $p$ \\
\midrule
1 (BIC-opt.) & 3 & 75{,}587 & 0\% & 33.7\% & 0.043 & 0.88 \\
2 & 15 & 75{,}625 & 0\% & 28.4\% & 0.031 & $< 0.001$ \\
3 & 8 & 75{,}660 & 0\% & 31.2\% & 0.038 & 0.001 \\
4 & 8 & 75{,}726 & 0\% & 29.1\% & 0.026 & $< 0.001$ \\
5 (econ.) & 7 & 75{,}805 & 90\% & 33.5\% & 0.019 & 0.018 \\
6 (econ.) & 3 & 75{,}906 & 100\% & 35.1\% & 0.022 & 0.17 \\
7 (econ.) & 6 & 76{,}137 & 92\% & 32.8\% & 0.026 & 0.28 \\
\bottomrule
\end{tabular}
\end{table}

This taxonomy reveals a fundamental challenge with purely statistical regime selection.
The highest log-likelihood fit (Cluster~1) achieves near-perfect statistical fit but
assigns crisis observations to non-crisis regimes, failing basic economic validity.
Cluster~5, with $\Delta\text{BIC} = 218$ versus Cluster~1, correctly identifies 90\% of
the 2008 GFC as Crisis states. The tension is unavoidable:
\begin{itemize}
\item \textbf{BIC criterion (Cluster~1):} Yields lowest fit criterion; fails to identify
the GFC as a distinct regime.
\item \textbf{Economic criterion ($\geq 50\%$ GFC detection, then highest LL):} Selects
Cluster~5; sacrifices $\Delta\text{BIC} = 218$ units to achieve crisis detection.
\end{itemize}

The OOS Elevated raw $F$-$p < 0.05$ in all 7 clusters, confirming that the
signal does not depend on cluster selection. However, per Section~\ref{sec:frozen_oos},
this robustness does not address prevalence-dependence (bootstrap $p = 0.153$).

The two-stage criterion (maximize GFC detection, then maximize likelihood within that constraint)
is pragmatic and economically motivated, but it is fundamentally post-hoc. We acknowledge this
limitation explicitly: had we relied on BIC alone, we would have obtained a fit inconsistent
with financial economics (zero GFC identification). Our approach prioritizes economic validity
over statistical parsimony, a choice that benefits robustness but cannot be justified as data-driven
in the strict sense.

\subsection{Scope, Interpretation, and Limitations}
\label{sec:limitations}

\textbf{Causal interpretation.}
As noted in Section~\ref{sec:granger_method}, all findings document predictive
precedence. The deleveraging hypothesis (Appendix~\ref{app:mechanism}) is plausible but
unverified; holdings-level 13F verification remains for future work.

\textbf{Economic magnitude and relevance.}
Effect sizes are small ($\Delta R^2 \approx 2\%$ pre-GFC) and do not generate
trading profits (Sharpe $= -0.07$; Appendix~\ref{app:trading}).
The contribution is not tradable alpha but \emph{risk model specification}:
multi-factor risk models (e.g., Barra, Axioma) that assume static cross-factor
covariance structure would have \emph{underestimated} tail risk during the
1998--2008 transition period when the HML$\to$SMB predictive link was active.
In a \$100B equity portfolio with 10\% HML/SMB exposure, a 2\% $\Delta R^2$
in cross-factor dynamics implies a $\sim$14\% increase in conditional SMB
variance during Normal-regime stress episodes (when the HML lag is at the
95th percentile), translating to $\sim$\$70M additional 1-day 99\% VaR.
The structural break finding implies that these risk model parameters
should be re-estimated after major market regime shifts---a practical diagnostic
even without direct trading implications.

\section{Conclusion}

\textbf{Primary finding: two-break structural decay.}
HML$\to$SMB Granger predictability is Bonferroni-significant in the
pre-crisis Normal regime ($p = 8.75 \times 10^{-9}$, $\Delta R^2 = 2.06\%$).
A data-driven sup-$F$ identifies June 1998 as the primary structural break
($p = 1.23 \times 10^{-13}$), with January 2008 confirming further decay
(Chow $p = 2.29 \times 10^{-6}$). The post-2008 coefficient has been
consistent with zero for 16 years (95\% CI $[-0.049, 0.073]$).
This is robust across all 7 HMM local optima clusters, multiple HAC bandwidth
specifications, lags 1--15 in Normal, and a trivariate specification conditioning on MKT-RF.

\textbf{Exploratory OOS evidence.}
Frozen OOS testing (HMM trained 1990--2012, applied 2013--2024 without refitting)
detects Elevated-regime re-emergence (HAC $p = 0.043$), but this
does not survive 30-pair Bonferroni correction. Prevalence sensitivity analysis
shows the signal is driven by regime expansion
(13.7\%$\to$33.7\%; bootstrap $p = 0.153$). The signal exists only at $K = 3$
and is bandwidth-sensitive (Andrews $p = 0.043$, NW default $p = 0.056$).
MOM$\to$SMB ($p = 0.010$) shows a similar pattern, suggesting the phenomenon
may extend beyond HML--SMB.

\textbf{Transfer entropy asymmetry.}
The reverse SMB$\to$HML channel is nonlinear and stronger than the
linear forward channel ($z = 5.37$ vs.\ $z = 2.45$ in Normal), suggesting
that linear connectedness measures systematically underestimate reverse
information flow.

\textbf{Implications.} Factor-timing models assuming regime-invariant cross-factor
relationships may misspecify dynamics during structural transitions. Future work should
pursue pre-registered replication on international factor data, direct verification
of deleveraging mechanisms via 13F holdings, and investigation of the nonlinear
reverse channel.

\section*{Funding}

This research received no specific grant from any funding agency in the public,
commercial, or not-for-profit sectors.

\section*{Data and Code Availability}

Code, analysis scripts, and the fixed seed list (50 seeds) for this paper are
available at an anonymous repository to be provided upon acceptance. The
financial data are publicly available from Kenneth French's data library at
\url{https://mba.tuck.dartmouth.edu/pages/faculty/ken.french/data_library.html}.

\emph{Reproducibility note:} Student-$t$ HMM estimation via EM is sensitive
to initialization (tolerance: $\Delta\text{LL} < 10^{-6}$; max iterations: 500;
all 50 seeds converged); independent researchers using different seeds may
obtain different local optima. We provide fixed seeds and tie-breaking rules
to enable exact replication of the primary fit (seed~28, selected as first
in sorted order among three seeds reaching the same full-sample log-likelihood).
The ``BIC-optimal'' label refers to $K = 3$ selection; the ``leakage-safe''
(LL-maximizing) label refers to seed selection within the $K = 3$ model,
which avoids using crisis-event alignment as a selection criterion.
The OOS Elevated result is significant in all 7 local optima clusters from
our 50-seed sweep. Permutation test seeds are fixed and documented in the
repository.

\appendix

\section{Hypothesized Economic Mechanism}
\label{app:mechanism}

We offer the following interpretation as a \emph{hypothesis}, not a verified mechanism:

\textbf{Normal-regime lag-1 link (pre-GFC, confirmatory):} During calmer periods,
cross-factor inventory management and short-horizon rebalancing may transmit
value shocks into size returns with a one-day delay.

\textbf{Stress-channel hypothesis (exploratory):} FF25 overlap is consistent with
a deleveraging channel in stress periods
($\rho_s = 0.35$, permutation $p = 0.046$). Primary Crisis-regime factor-level
Granger tests remain null in the leakage-safe fit.

\textbf{Caveat: dual-mechanism problem.} The Normal-regime ``inventory management''
story differs from the Elevated/Crisis ``deleveraging cascade'' hypothesis.
That different regimes require different mechanisms suggests either (a)~multiple
channels operate at different volatility levels, or (b)~the statistical pattern
reflects non-economic regularities. This ambiguity should temper confidence in
any single economic interpretation.

\textbf{Alternative explanations we cannot rule out:}
\begin{itemize}
    \item A common driver (e.g., funding liquidity) affects both factors with different lags
    \item Information about distress propagates through markets, reaching Value-sensitive
    investors before Size-sensitive investors
    \item Statistical artifact of regime-conditional estimation or $K = 3$ specification
\end{itemize}

Greenwood and Thesmar~\cite{greenwood2011stock} document that stocks with
concentrated institutional ownership exhibit correlated fragility during
liquidation events---consistent with our hypothesized deleveraging channel.
We provide exploratory portfolio-level evidence for this overlap in
Section~\ref{sec:overlap_evidence}.

\section{Portfolio Overlap Evidence}
\label{app:overlap}

The deleveraging hypothesis predicts that the HML$\to$SMB Granger link should
be strongest for portfolios that load on \emph{both} factors simultaneously.
The Small\-/HighBM portfolio sits in both the SMB long leg (small stocks) and
the HML long leg (high book-to-market), making it the primary candidate
transmission channel.

We test this prediction using the Fama-French 25 Size$\times$BM portfolios
(daily, 1990--2024). For each of the 25 portfolios, we run a Granger test
of HML$\to$Portfolio at lag 1 (BIC-selected in Table~\ref{tab:main}) within
the Crisis regime.

\textbf{Spatial gradient test (exploratory).} We define an ``overlap score'' for each
portfolio as $(4 - \text{size quintile}) + \text{BM quintile}$, ranging from
0 (Big/LowBM, no overlap) to 8 (Small/HighBM, maximal overlap). The
confirmatory test is a Spearman rank correlation between overlap score and
$-\log_{10}(p)$ of the Granger test across 25 portfolios, with significance
assessed by spatial permutation (10{,}000 shuffles of grid assignments;
time series are never touched).

\textbf{Result:} $\rho_s = 0.35$, permutation $p = 0.046$
(Figure~\ref{fig:ff25}). Granger significance concentrates in small-cap
portfolios: all five Small quintile portfolios are nominally significant
($p < 0.01$), with three surviving Bonferroni correction ($\alpha/25 = 0.002$):
S/L ($p = 2.5 \times 10^{-6}$, $\Delta R^2 = 1.47\%$),
S/3 ($p = 0.0014$), and S/H ($p \approx 0$, $\Delta R^2 = 1.79\%$).
One Big-quintile portfolio is nominally significant ($p < 0.01$), but none
survive Bonferroni. This spatial pattern is inconsistent
with a uniform factor-level effect and supports the overlap channel.

\textbf{Overlap fraction.} At lag 15,
\[
\frac{\Delta R^2(\text{HML}\to\text{S/H})}{\Delta R^2(\text{HML}\to\text{SMB})}=0.392
\]
Small\-/HighBM alone accounts for 39\% of the HML$\to$SMB incremental $R^2$.

\begin{figure}[t]
\centering
\includegraphics[width=\columnwidth]{figures/ff25_overlap_granger_heatmap.pdf}
\caption{FF 25 Portfolio Overlap Analysis. Panel A: $-\log_{10}(p)$ for\\
HML$\to$Portfolio Granger tests (Crisis, lag=1). Significance concentrates in
Small quintile; $^{**}$Bonferroni-significant. Panel B: Spatial gradient---overlap
score vs.\ Granger significance ($\rho_s = 0.35$, permutation $p = 0.046$).}
\Description{Two-panel heatmap showing FF25 portfolio Granger significance by size and BM quintile, with spatial gradient plot.}
\label{fig:ff25}
\end{figure}

\section{Event-Based Consistency Checks}
\label{app:events}

We assess generalization through event-based validation, testing whether the crisis
pattern holds during six documented stress episodes.

\begin{table}[t]
\centering
\caption{Event-Based Consistency Checks. Expected pattern: HML $\to$ SMB significant during
Crisis. Result codes: \checkmark = expected pattern holds ($p < 0.10$ for HML$\to$SMB,
$p > 0.10$ for reverse); Dir.\ = correct direction but insufficient power;
$\times$ = opposite pattern.}
\scriptsize
\begin{tabular}{@{}lcccc@{}}
\toprule
Event & Days & $p$(H$\to$S) & $p$(S$\to$H) & Result \\
\midrule
2008 Financial & 189 & 0.118 & 0.648 & Dir. \\
2011 EU Debt & 85 & 0.090 & 0.321 & \checkmark \\
2015 China & 64 & 0.419 & 0.459 & Dir. \\
2018 Vol Shock & 29 & 0.236 & 0.179 & $\times$ \\
2020 COVID & 83 & 0.062 & 0.134 & \checkmark \\
2022 Rate Hikes & 188 & 0.626 & 0.112 & $\times$ \\
\bottomrule
\end{tabular}
\end{table}

\textbf{Summary (exploratory, low-power):} The expected HML$\to$SMB pattern
holds in 2 of 6 events (2011, 2020), with 4 of 6 directionally correct
(binomial $p \approx 0.11$, not significant). Two exceptions show reversed dynamics:
the 2018 Vol Shock ($n = 29$) is too short for reliable estimation, while the
2022 Rate Hikes involved a value-growth rotation driven by monetary policy---a
mechanism distinct from the liquidity-driven deleveraging channel, where
interest-rate sensitivity may reverse the usual HML$\to$SMB ordering.
These reversals temper the mechanistic interpretation but do not undermine
the aggregate pooled result, which averages across heterogeneous event types.

\section{Markov-Switching VAR Comparison}

A single-stage MS-VAR~\cite{krolzig1997markov} ($K = 3$, 9 lags, switching variance)
confirms HML lags are informative ($\text{LR} = 114.95$, $df = 27$, $p = 8.0 \times 10^{-13}$),
though BIC favors the restricted model ($\Delta\text{BIC} = -130$) due to the 27 additional
parameters. The HML--SMB link exists within a principled single-stage framework
but its magnitude is too small to justify the model complexity.

\section{Trading and Risk Monitoring}
\label{app:trading}
\label{app:var}

\textbf{Trading backtest.}
A proof-of-concept strategy---long SMB when 9-day cumulative HML is positive
during Crisis, short otherwise---yields Sharpe $= -0.07$ (vs.\ buy-and-hold
$+0.06$) over 1995--2024. This intentionally simple rule (no optimization,
no transaction costs) confirms that \emph{statistical predictability}
$\neq$ \emph{economic predictability}; it does not rule out
profitability under more sophisticated implementations.

\textbf{VaR application.}
Among five VaR specifications, regime-conditional models pass conditional coverage tests but exhibit high false-alarm rates (93.2\%), suggesting simpler GARCH alternatives may be preferable.

\section{Extended Robustness Analyses}
\label{app:robustness}

\textbf{Subsample stability.} Splitting at 2008: Crisis HML$\to$SMB
pre-2008 ($n = 659$, $p = 0.335$) and post-2008 ($n = 322$, $p = 0.599$)
are both non-significant, so we avoid claims of subsample-confirmed Crisis effects.

\textbf{Alternative regime methods.} Threshold-based volatility regimes
(realized volatility $>$ 90th percentile = Crisis) also detect the pattern
in sensitivity checks, but we treat this as exploratory given the primary-fit
null in Crisis at lag 1.

\textbf{Practitioner baselines (exploratory).} Standard unconditional tools detect the HML$\to$SMB
asymmetry but miss the architectural distinction. Rolling 250-day Granger tests
find HML$\to$SMB significant in 36.2\% of windows vs.\ 9.6\% for the
reverse---with concentration in Normal (47.1\% vs.\ 13.7\% reverse),
lower rates in Elevated (21.3\%), and mixed Crisis windows (27.7\%).
FEVD spillover is regime-insensitive in direction,
failing to detect the dormant Elevated regime or the Normal/Crisis directionality
shift. Regime conditioning does not merely sharpen detection---it reveals
qualitatively different predictive architectures that rolling methods average away.

\textbf{Weekly aggregation.} Using weekly returns, the HML$\to$SMB crisis pattern
remains non-significant ($p = 0.114$, $n = 221$), consistent with weak
evidence outside Normal.

\textbf{Regime transitions (exploratory).} Analyzing 60-day windows around Normal$\to$Crisis
transitions, the HML$\to$SMB relationship strengthens upon crisis entry
(median $p$: $0.058 \to 0.021$), suggesting discrete intensification of the
predictive link.

\textbf{ML diagnostic protocol details} (supporting Section~\ref{sec:ml_validation}).
We apply four model classes---Linear (OLS), RF (100 trees), MLP (64-32),
LSTM (32 hidden, 9-step)---to characterize
the \emph{complexity boundary} of the predictive relationship.
Each predicts SMB from its own lags (restricted) vs.\ both SMB and
HML lags (unrestricted), with permutation testing
(200 shuffles for Linear/RF/MLP, 100 for LSTM; at 200 permutations,
$p$-value SE $\approx 0.007$ at $p = 0.01$) and expanding-window CV.

Table~\ref{tab:neural} and Figure~\ref{fig:complexity} present the results.
Linear improvements are modest (Normal 0.86\%, Elevated 0.92\%, Crisis 0.43\%).
Under the primary fit (seed~28), nonlinear permutation tests are non-significant
($p > 0.16$ throughout); however, under seed~42, MLP is marginally significant
($p = 0.015$), indicating fit-dependence.\footnote{The LSTM was estimated under
a separate HMM initialization; $p$-values are approximate due to the 100-permutation limit.}
This complexity characterization supports a conservative conclusion:
incremental predictive content in the HML$\to$SMB direction is captured by
linear terms, though TE reveals residual nonlinear structure in the reverse
direction (Table~\ref{tab:te}).

\begin{figure}[t]
\centering
\includegraphics[width=\columnwidth]{figures/complexity_spectrum.pdf}
\caption{Four-model diagnostic protocol: MSE improvement (\%) from including
HML lags, by regime and model class. Only linear models (blue) achieve
significant improvement in Normal and Elevated; all nonlinear methods (RF, MLP, LSTM)
fail to improve prediction.}
\Description{Bar chart comparing MSE improvement across four model classes and three regimes, showing only linear models achieve significant improvement.}
\label{fig:complexity}
\end{figure}

\begin{table}[t]
\centering
\caption[Four-model diagnostics]{Four-Model Diagnostic Protocol: HML$\to$SMB by Regime.
Linear column shows MSE improvement (\%); others show permutation $p$-values.}
\label{tab:neural}
\small
\begin{tabular}{lccccc}
\toprule
Regime & $n$ & Linear & RF $p$ & MLP $p$ & LSTM $p$ \\
\midrule
Normal & 4,496 & 0.86\%$^{**}$ & 0.69 & 0.20 & 0.63$^\dagger$ \\
Elevated & 2,792 & 0.92\%$^{**}$ & 1.00 & 1.00 & 0.93$^\dagger$ \\
\textbf{Crisis} & \textbf{1,017} & \textbf{0.43\%} & 0.13 & 0.96 & 0.83$^\dagger$ \\
\bottomrule
\multicolumn{6}{l}{\footnotesize $^{**}p < 0.01$. $^\dagger$LSTM: separate HMM init., 100 permutations.} \\
\end{tabular}
\end{table}

\textbf{Transfer entropy details} (supporting Section~\ref{sec:ml_validation}).
Table~\ref{tab:te} reports multi-lag TE~\cite{schreiber2000measuring,barnett2009granger}
(Frenzel--Pompe kNN, $k=5$, 200 permutations).
The $k=5$ choice follows standard recommendations for kNN entropy estimation;
sensitivity checks at $k \in \{3,4,5,6,7\}$ yield qualitatively identical
directional asymmetry (forward $z \in [2.1, 2.8]$, reverse $z \in [4.9, 5.6]$
in Normal), confirming the finding is not an artifact of the kNN parameter.
The directional asymmetry---linear forward HML$\to$SMB but nonlinear
reverse SMB$\to$HML---is discussed in Section~\ref{sec:ml_validation};
the ``predominantly linear'' characterization applies strictly to the
\emph{forward} direction under the primary fit.

\begin{table}[t]
\centering
\caption{Transfer Entropy: Multi-Lag (1--9) by Regime.
$z$-scores relative to permutation null (200 shuffles).}
\label{tab:te}
\scriptsize
\begin{tabular}{@{}lccccc@{}}
\toprule
Regime & $n$ & HML$\to$SMB $z$ & $p$ & SMB$\to$HML $z$ & $p$ \\
\midrule
Normal & 4,496 & $2.45^{**}$ & 0.007 & $5.37^{***}$ & $< 10^{-6}$ \\
Elevated & 2,792 & $2.41^{**}$ & 0.008 & $1.65^{*}$ & 0.049 \\
Crisis & 1,017 & $1.01$ & 0.157 & $1.22$ & 0.111 \\
\bottomrule
\multicolumn{6}{l}{\footnotesize $^{***}p < 10^{-6}$, $^{**}p < 0.01$, $^{*}p < 0.05$.} \\
\multicolumn{6}{l}{\footnotesize $n$ = clean observations after boundary exclusion.} \\
\end{tabular}
\end{table}

\textbf{Filtered vs.\ smoothed.} Overall agreement: 95.9\%.
Crisis HML$\to$SMB under filtered probabilities: $p = 0.601$ (vs.\ $0.695$ smoothed).
Qualitatively consistent across classification methods.

\section{Local Optima Taxonomy and Boundary Exclusion Details}
\label{app:optima}

This section provides technical details supplementing the main-text discussion in Section~\ref{sec:local_optima},
which presents the taxonomy and its economic implications in condensed form.

\textbf{Local optima taxonomy.}
HMM estimation via EM is sensitive to initialization. Our 50-seed multistart
reveals 7 distinct local optima clusters:
Cluster~1 (3 seeds, best BIC~$= 75{,}587$, 0\% GFC detection, Crisis $p = 0.88$);
Cluster~2 (15 seeds, BIC~$= 75{,}625$, 0\% GFC, Crisis $p < 0.001$);
Cluster~3 (8 seeds, BIC~$= 75{,}660$, 0\% GFC, Crisis $p = 0.001$);
Cluster~4 (8 seeds, BIC~$= 75{,}726$, 0\% GFC, Crisis $p < 0.001$);
Cluster~5 (7 seeds, BIC~$= 75{,}805$, 90\% GFC, Crisis $p = 0.018$);
Cluster~6 (3 seeds, BIC~$= 75{,}906$, 100\% GFC, Crisis $p = 0.17$);
Cluster~7 (6 seeds, BIC~$= 76{,}137$, 92\% GFC, Crisis $p = 0.28$).
Only clusters with $\geq$50\% GFC detection (5--7) are economically valid for
crisis-regime inference; among these, Cluster~5 has the best fit.
The highest log-likelihood cluster (seeds 28, 20, 6; BIC~$= 75{,}587$)
assigns 0\% of 2008 crisis days to the Crisis regime ($p = 0.88$),
while a lower-likelihood cluster (seeds 21, 9, 0, 29, 46, 36, 11;
BIC~$= 75{,}805$, $\Delta\text{BIC} = 218$) assigns 90\% of 2008 to Crisis
($p = 0.018$). Readers who prefer pure BIC selection should interpret the
Crisis-regime Granger results as seed-dependent; the Elevated-regime OOS results
(primary seed 28) are robust across all 7 clusters ($p < 0.05$ in 7/7),
though this robustness is conditional on the full OOS Elevated sample;
bootstrap prevalence reweighting yields median $p = 0.153$
(Section~\ref{sec:frozen_oos}).

\textbf{Boundary exclusion.}
In-sample: 0.67\% (59/8{,}817) excluded; OOS: 7.4\% (224/3{,}020).
Median inter-regime gap: 3 calendar days; exclusion rates are not
systematically biased toward any regime. This conservative approach
may underestimate transition effects but ensures clean within-regime estimation.

\bibliography{references}

\end{document}
